\cotefile{91}
\scanpage{59}
% archive foliation: 58
\unsure{comme pour}
\[
T\longrightarrow f_*^{\mathrm{ét}.}(T)
\]
qui est \unsure{tel} que l'hom canonique
\[
\underline{\mathrm{Hom}}_S(\overbrace{f_*^{\mathrm{ét}}(T)}^{\inserted{\overline T}},S')\longrightarrow
\underline{\mathrm{Hom}}_S(T,S')
\]
est un isom. pour tt $S'$. (\unsure{raison} --). Il en résulte que pour tt \deleted{\unsure{$S'\to S$}} changement de base
\deleted{$\underline{\mathrm{Hom}}_{\overline{k(s)}}(\overline T_{\bar s},S'_{\bar s})\to$}
(\unsure{p. ex.} $\bar s=\operatorname{Spec}(\overline{k(s)})$)
la \deleted{\unsure{a}} relation reste vraie; faisant $S'=S\times I$ ($I$ \unsure{un fini}) on en conclut que
$T\to f_*^{\mathrm{ét}}(T)$ est surjectif et \inserted{(si $T/S$ est fini)} géométriquement injectif
(\unsure{regarder} \deleted{\unsure{un}} \inserted{un ensemble} \unsure{ch.} pour tt $s\in S$).

Inversement, si $T\to\overline T$ est un morphisme surjectif et géom. injectif de $T$ (fini \inserted{et plat} sur $S$) sur $\overline T$ \unsure{sch.} étale de $S$, alors on a
\[
\underline{\mathrm{Hom}}_S(\overline T,S')=\mathrm{Hom}_S(T,S')
\]
pour tout $S'$ non ramifié sur $S$ (\unsure{enfin} $\unsure{f}$ est \unsure{comme} \unsure{le} $=$ \unsure{dessus} reste vrai par \unsure{chgt.} de la base\inserted{$(S)$}),
\unsure{encore} \deleted{\unsure{$\mathrm{H}$}}
$\underline{\mathrm{Hom}}_S(\overline T,S')\simeq\underline{\mathrm{Hom}}_S(T,S')$,
\unsure{isom.} \unsure{T---}.
